\documentclass[landscape, 8pt]{extarticle}

\usepackage{../../preamble}
\usepackage[separate]{../../rss/thmboxes_v5}
\usepackage{bm}

%I based these notes off the course notes with 0 knowledge as I attended very few of the lectures and did not engage with assignments/workshops (I broke my thumb in Week 1 of this semester). Use at your own risk!


\begin{document}
	
	\fontsize{5pt}{6pt} \selectfont
	
	%Reduce horizontal padding in 'tabular' environment, make rows taller for legibility
	\setlength\tabcolsep{4.5pt}
	\setlength\extrarowheight{2pt}
	
	\setlength{\abovedisplayskip}{0pt}
	\setlength{\belowdisplayskip}{0pt}
	\setlength{\abovedisplayshortskip}{0pt}
	\setlength{\belowdisplayshortskip}{0pt}
	
	\begin{multicols*}{4}
		\raggedcolumns 
		%Original template by Leon, modified by Louis, content by Louis
		\vspace{-5pt}
		\subsection*{CMM Exam Sheet 2025/26 - Non ii, non vidi, non vicī}
		\vspace{-5pt}
		
		
		\begin{rem}[Dimension]{rem:dim}{ }
			Dimension is very useful when sanity-checking answers. The basic dimensions we deal with are length $(L)$, time $(T)$, and mass $(M)$. For example, $\bm{x}\in \mathbb{R}^3$ has a dimension of length, often written $[\bm{x}] = L$. Similarly the time-derivative has dimension of inverse time, so if $Q(t)$ has dimension $[Q]$, its time-derivative has dimension $[\dot{Q}] = [Q]T^{-1}$. Velocity and acceleration have dimensions $[\bm{\dot{x}}] = LT^{-1}$ and $[\bm{\ddot{x}}] = LT^{-2}$. Dimension is multiplicative - $[Q_1Q_2] = [Q_1][Q_2]$. For example, $[\bm{F}] = [m][\bm{\ddot{x}}] = MLT^{-2}$.
			
		\end{rem}
		
		
		\begin{dfn}[The Newtonian Universe]{dfn:newtonuni}{ }
			\begin{itemize}
				\item An \textbf{event} is some $(t, \bm{a})\in \mathbb{R} \times \mathbb{R}^3$ labelling a fixed $t$ in time and $\bm{a}$ in space.
				\item Two events $(t_0, \bm{a}), (t_1, \bm{b})$ are \textbf{simultaneous} iff $t_0 = t_1$. The distance between simultaneous events is the norm $\left|\bm{a} - \bm{b}\right|$
				\item The building blocks of our systems are \textbf{particles}, objects whose position in the universe are given at any time by one event
				\item A \textbf{worldline} is a particle's trajectory through the universe. It is a graph $\bm{x} : \mathbb{R} \to \mathbb{R}^3, t \mapsto\bm{x}(t)$ of form $(t, \bm{x}(t))\in \mathbb{R} \times \mathbb{R}^3$.
				\item \textbf{Velocity} and \textbf{acceleration} are $\bm{v} := \bm{\dot{x}}$ and $\bm{\dot{v}} := \bm{\ddot{x}}$ both w.r.t. $t$
				\item The \textbf{configuration space} of a \textbf{many-body system of $n$ particles} is the location of each particle in space, so is $\mathbb{R}^3 \times \mathbb{R}^3 \times \dots \times \mathbb{R}^3 = \mathbb{R}^{3n} = \mathbb{R}^N$.
				\item For the worldline of the $a^{th}$ particle denoted $\bm{x}_a(t): \mathbb{R} \to \mathbb{R}^3$, the collection of the worldlines defines a curve $\bm{x}(t) = (\bm{x}_1(t), \bm{x}_2(t), ..., \bm{x}_n(t)): \mathbb{R} \to \mathbb{R}^N$ in the configuration space.
			\end{itemize}
			
		\end{dfn}
		
		
		\begin{dfn}[Newton's Principle of Determinacy]{dfn:newtdetermprinc}{ }
			The motion of any mechanical system is uniquely determined by its initial state, the set of all of its constituent positions and velocities at a given instant in time. 
			
		\end{dfn}
		
		
		\begin{mdfn}[Newton's Equation]{dfn:newtoneq}{ }
			\textbf{Newton's Equation} is the second-order ODE $\bm{\ddot{x}}(t) = \bm{\Phi}(\bm{x}, \bm{\dot{x}}, t)$ for $\bm{\Phi} : \mathbb{R}^n \times \mathbb{R}^N \times \mathbb{R}$. It is the mathematical embodiment of Newton's principle because the constants of integration in a general solution to the ODE are fixed by the initial state.\\
			As a system of first-order ODEs: Let $\bm{v} : \mathbb{R} \to \mathbb{R}^N$, then the system $\binom{\bm{\dot{x}}(t)}{\bm{\dot{v}}(t)} = \binom{\bm{v}(t)}{\bm{\Phi}(\bm{x}, \bm{v}, t)}$ has the same solution space.
			
		\end{mdfn}
		
		
		\begin{mthm}[Newton's Laws Of Motion]{mthm:newtlaws}{ }
			\begin{enumerate}[label=(\arabic*)]
				\item A particle (or system of particles) is \textbf{freely-moving} if $\bm{\Phi}(\bm{x}, \bm{\dot{x}}, t) = 0$
				\begin{itemize}
					\item A freely moving system moves on a straight worldline at constant $\bm{v}$
				\end{itemize}
				\item A \textbf{force} is a map $\bm{F} : \mathbb{R}^3 \to \mathbb{R}$ affecting a particle's $\bm{v}$. For a particle of inertial mass $m$, N's E in a force field is $\bm{F}(\bm{x}) = m \bm{\ddot{x}}$ or $\bm{\Phi}(\bm{x}, \bm{\dot{x}}, t) = \frac{\bm{F}(\bm{x})}{m}$.
				
			\end{enumerate}
			
		\end{mthm}
		
		
		\begin{dfn}[Affine Space]{dfn:affnspc}{ }
			
			
		\end{dfn}
		
		
		\begin{dfn}[Galilean Transformations]{dfn:galiltransforms}{ }
			A \textbf{Galilean transformation} is an affine transformation which preserves time intervals and the distance between simultaneous events
			\begin{itemize}
				\item A Galilean transformation is given by a linear transformation of the form $\begin{psmallmatrix*}
					1&\bm{0}&s\\
					\bm{v}&R&\bm{s}\\
					0&\bm{0}&1
				\end{psmallmatrix*}$ where $\bm{s, v}\in \mathbb{R}^3, s\in \mathbb{R},$ and $R: \mathbb{R}^3 \to \mathbb{R}^3$ is an orthogonal transformation.
			\end{itemize}
			
		\end{dfn}
		
		
		
	\end{multicols*}
	
	
\end{document}
